\subsection{What a small cycle supplies at the common entrance}

Use the fixed-field orientation $\ell\ge0$ of
Section~\ref{rw:geometry}, with the reflection-symmetric thresholds
of Section~\ref{rw:counting}. The coefficient divisions are those
of the original signed auxiliary maps.
The source representative counts the small cycles and determines the
ratio of the two homogeneous generators. To obtain the entrance sign
of $K$, we then evaluate the regular comparison in the original central
height. This change of representation is necessary because a unit
relating closing functions need not preserve their derivative zeros.

\begin{proposition}\label{rw:small-entrance}
After a common neighborhood restriction, suppose $\rho>0$ and
\eqref{rw:counting:forcing} holds. There is at most one counted
cycle with $p<p_\sharp/4$, and it is simple. A counted cycle with
$p<p_\sharp/8$ implies
\begin{equation}\label{rw:upper:small-ratio}
 \tau_1\sigma_1<0,\qquad
 0<-\tau_1/\sigma_1\le p_\sharp^2/300.
\end{equation}
The actual regular comparison at $p_0=p_\sharp/4$ then satisfies
$\operatorname{sign}K(p_0)=-\operatorname{sign}\sigma_1$.
Thus it supplies the precise signed entrance of hypothesis
\textup{(III)} in Section~\ref{rw:counting}.
\end{proposition}
\begin{proof}
The endpoint action bound of Section~\ref{rw:finite:geometry},
with the fixed primary normalization, places every parameter-dominant
cycle beyond $p>\eta/8\ge p_\sharp$; the same bound is recorded
in the physical cover of Section~11. Hence the small cycles here
belong to the two initial energy classes, and are zeros of the one
function constructed in Proposition~\ref{rw:initial-joint-bound}.
In the source class its outgoing action satisfies
\[
 p^2=2e_+\epsilon h,\qquad |e_+-1|<1/1024.
\]
For a cycle with $p<p_\sharp/4$ this gives
$\epsilon h<p_\sharp^2/31$. On the compact energy class the
same inequality follows from $h\le H_*$ after the scale cutoff
is reduced. Let $J_0$ be the initial-energy interval up to that
upper bound; it contains all these cycles, with their multiplicities.
Its lower energy bound is the fixed $h(z_L)>H_{\rm ECT}+2$.
We do not use any ECT assertion near energy zero.

Write the actual source representative as
\[
 F_{\rm src}=\tau_0 f_1+\sigma_0 f_2+f_3,
 \qquad Q=f_2/f_1,\quad g=f_3/f_1.
\]
The compact overlap and the source derivative estimates in Section~6
give on this single initial-energy interval
\[
 Q\ge q_*>0,\qquad Q'\ge q_0>0,\qquad Q\le C_Qh.
\]
On the compact interval these inequalities follow from the first two
ordered Wronskians and fixed coefficient convergence; on the source
tail they follow from $Q=hC/A$, $Q_h>1/5$. The fixed overlap joins
the same two columns, so their signs can be tested on the union.

On $J_0$ we need the line-breaking contribution and its ordinary
energy derivative to be exponentially smaller than the homogeneous
coefficient norm. The signed compact contribution is $\ell$ times a
bounded coefficient. On the source part it is the exact logarithmic
forcing divided by $h^{3/2}A$.
Since $\widehat\ell=\ell/s$ and $s^{-1}=e^{h+1-s}\le e^{h+1}$,
\eqref{rw:counting:forcing} and the source logarithmic estimates
yield
\begin{equation}\label{rw:upper:small-forcing}
 |g|+|g'|\le C\exp\left(1-\frac{11p_\sharp^2}{620\epsilon}\right)
                      (|\tau_0|+|\sigma_0|)
       =:\vartheta_0(|\tau_0|+|\sigma_0|).
\end{equation}
Here $1/20-1/31=11/620$. The same bound applies to the fixed cutoff
terms on the compact overlap. The source logarithms are evaluated at
$\widehat\ell=o(1)$ on this smaller interval, so their denominators
have fixed margins. The derivative is the actual fixed-coefficient
energy derivative, not an independent source derivative.

If both homogeneous generators vanish, forcing gives $\ell=0$,
and the center identity contributes no isolated zero. Otherwise
\[
 F_{\rm src}/f_1=\tau_0+\sigma_0Q+g,
 \qquad (F_{\rm src}/f_1)'=\sigma_0Q'+g'.
\]
If $|\sigma_0|q_0>\vartheta_0(|\tau_0|+|\sigma_0|)$,
this derivative has a strict sign on all of $J_0$. Every zero is
therefore simple and there is at most one. In the complementary
case $|\sigma_0|\le2\vartheta_0|\tau_0|/q_0$ at small scale,
and
\[
 |F_{\rm src}/f_1|\ge |\tau_0|
 \left(1-\frac{2\vartheta_0 C_Qp_\sharp^2}{31q_0\epsilon}
                         -2\vartheta_0\right)>0.
\]
The exponential absorbs $\epsilon^{-1}$, so there is no zero.
This proves uniqueness and simplicity on the whole small-cycle set.

At the smaller cycle, with energy $h_y$, the same equation gives
\[
 -\frac{\tau_0}{\sigma_0}
 =Q(h_y)+g(h_y)/\sigma_0,
 \qquad |g/\sigma_0|\le\vartheta_0(1+|\tau_0/\sigma_0|).
\]
For small $\vartheta_0$, absorb the ratio on the right. The positive
lower bound $Q\ge q_*$ proves the strict negative sign of
$\tau_0/\sigma_0$. Moreover $\epsilon h_y<p_\sharp^2/127$.
The independent source constants give $C/A\to2/5$. Choose a fixed
$H_1\ge H_*+1$ so that its scale-tail error is below $1/100$,
then choose $p_\sharp$ small so that the coefficient error on this
energy range is below $1/100$. For $h\ge H_1$, $Q\le0.42h$;
below it $Q\le C_QH_1$. Therefore
\[
 \epsilon Q(h_y)\le
 \max(0.42p_\sharp^2/127,\epsilon C_QH_1)
                    \le p_\sharp^2/302
\]
after reducing the primary scale. The strict inequality
$0.42/127<1/302$ and the exponential error leave the margin to
$1/300$. Since $\tau_1/\sigma_1=\epsilon\tau_0/\sigma_0$,
this proves \eqref{rw:upper:small-ratio}.

The generator ratio now determines the entrance sign of the original
central comparison.
The physical-fibre construction in Section~11 includes the regular
point $p_0=p_\sharp/4$. The two regular columns in
Lemma~\ref{rw:regular-columns} satisfy
$C/A=p^2/5+O(p^3)$, uniformly on the normalized compact coefficient
set. Thus at this fixed point $C/A\ge p_\sharp^2/85$ after
choosing $p_\sharp$ small. Let
$w=(\tau_1,\sigma_1)/(|\tau_1|+|\sigma_1|)$.
Equation \eqref{rw:upper:small-ratio} makes
$w_1+w_2C/A$ have sign $\sigma_1$ and modulus at least
$p_\sharp^2/200$. The exact regular representation, divided by
this generator norm, is
\[
 K(p_0)/(|\tau_1|+|\sigma_1|)
          =-A(p_0)(w_1+w_2C(p_0)/A(p_0))+e(p_0).
\]
The positive $A(p_0)$ has a fixed margin, and regular coefficient
convergence together with the fixed-field forcing makes $e(p_0)$
smaller than half the preceding margin. Consequently $K(p_0)$ has
sign $-\sigma_1$, as required. All choices preceded the fixed field
and its cycle. Reflection transfers this statement with its actual
outgoing label and the identity \eqref{rw:counting:reflection}.
\end{proof}
